Deep image classifiers routinely reach high top-1 accuracy on benchmarks
such as ImageNet, but high accuracy does not by itself tell us \emph{why} a
model reaches a given prediction. This matters most at the margin: when a
model confuses two visually similar classes, understanding what the model
was actually looking at can reveal whether the error is a reasonable mistake
(the model attended to the right region but a genuinely ambiguous or
occluded cue) or a symptom of a shortcut (the model ignored the
discriminating feature entirely).

SHAP~\cite{lundberg2017unified} is a model-agnostic explanation method that
assigns each input feature -- in our case, image regions -- a contribution
score toward a specific output. Applied to image classification, SHAP
attribution maps let us visually inspect which regions of an image pushed
the model toward its predicted label.

This paper asks a narrow, concrete question: \emph{when an image classifier
confuses two visually similar classes, does SHAP show that the model
attended to the feature a human would use to tell them apart (ear shape,
coat pattern, beak shape), or does it attend elsewhere?} We answer this by
building a small but systematic test set of confusable ImageNet class pairs,
running a pretrained ConvNeXt-Tiny classifier, and inspecting SHAP
attribution maps for both correct and incorrect predictions.
